\documentclass[10pt]{article}
\usepackage[english]{babel}
\newcommand{\nombre}{Luis Valverde Ramos}
\newcommand{\pie}{Curriculum vitae, November 2026}
% Diseño compartido de cv-es.tex y cv-en.tex. Compila con pdfLaTeX (también
% en Overleaf): Archivo es la misma letra de la web y el verde es su acento.
% Las referencias van en privado/referencias-<idioma>.tex, fuera del repo, y
% solo entran si se compila con \def\conreferencias{} (build.py --privado).

\usepackage[a4paper,left=1.8cm,right=1.8cm,top=1.4cm,bottom=1.6cm,footskip=0.8cm]{geometry}
\usepackage[T1]{fontenc}
\usepackage[lining,tabular,mainfont,semibold]{Archivo}
\usepackage{microtype}
\usepackage{xcolor}
\usepackage{enumitem}
\usepackage{needspace}
\usepackage{fancyhdr}
\usepackage{lastpage}
\usepackage[hidelinks]{hyperref}

\definecolor{acento}{HTML}{2E6B4F}
\definecolor{tinta}{HTML}{191E1B}
\definecolor{gris}{HTML}{525A54}
\definecolor{linea}{HTML}{C2C9BF}

\setlength{\parindent}{0pt}
\linespread{1.05}
\frenchspacing
\brokenpenalty=10000 % ninguna palabra partida entre dos páginas
\AtBeginDocument{\color{tinta}}

% Columna de fechas a la izquierda y texto a la derecha
\newlength{\fechas}\setlength{\fechas}{2.45cm}
\newlength{\separa}\setlength{\separa}{0.4cm}

% Pie: nombre, fecha de actualización y página
\pagestyle{fancy}
\fancyhf{}
\renewcommand{\headrulewidth}{0pt}
\fancyfoot[C]{\scriptsize\color{gris}\nombre\ · \pie\ · \thepage/\pageref*{LastPage}}

% Cabecera: nombre, una línea de presentación y los contactos
\newcommand{\cabecera}[2]{%
  {\fontsize{23}{27}\selectfont\bfseries\nombre}\par\vspace{4pt}
  {\color{gris}#1}\par\vspace{3pt}
  {\small #2}\par\vspace{7pt}
  {\color{acento}\hrule height 1.2pt}\vspace{2pt}}
\newcommand{\enlace}[2]{\href{#1}{\color{acento}#2}}

% Sección: título en versalitas verdes con una raya fina; nunca sola al pie
\newcommand{\seccion}[1]{%
  \par\Needspace{5\baselineskip}\vspace{9pt}%
  {\color{acento}\bfseries\footnotesize\textls[110]{\MakeUppercase{#1}}}\par
  \vspace{2pt}{\color{linea}\hrule height 0.5pt}\vspace{6pt}}

% Lista de entradas con la fecha en la columna izquierda
% \leavevmode antes del color, para que la fecha comparta la línea base del
% cargo; \smash, para que una fecha de dos líneas no empuje el texto
\newcommand{\fecha}[1]{\smash{\parbox[t]{\fechas}{\raggedright\small\leavevmode\color{gris}#1}}}
\newlist{entradas}{description}{1}
\setlist[entradas]{leftmargin=\dimexpr\fechas+\separa\relax, labelwidth=\fechas,
  labelsep=\separa, align=left, font=\normalfont, itemsep=5pt, topsep=0pt,
  parsep=1pt, partopsep=0pt}
% Cada entrada empieza en una página con sitio para su título y dos líneas más
\newcommand{\entrada}[2][]{\Needspace{3\baselineskip}\item[\fecha{#1}]#2}

% Cargo e institución, y un detalle opcional en gris debajo
\newcommand{\cargo}[2]{{\bfseries #1}, #2\par}
\newcommand{\detalle}[1]{{\small\color{gris}#1}\par}

% Viñetas dentro de una entrada
\newlist{puntos}{itemize}{1}
\setlist[puntos]{label={\color{acento}\textbullet}, leftmargin=1.1em,
  labelsep=0.45em, itemsep=1.5pt, topsep=2pt, parsep=0pt, before=\small}


\begin{document}

\cabecera{Political scientist (PUCP) · Lima, Peru}{%
  \enlace{mailto:luis.valverde@pucp.edu.pe}{luis.valverde@pucp.edu.pe} \ · \
  \enlace{https://luiseduardovalverderamos.github.io/}{luiseduardovalverderamos.github.io} \ · \
  \enlace{https://www.linkedin.com/in/luisvalverderamos/}{linkedin.com/in/luisvalverderamos}}

\seccion{Profile}
Political scientist trained at the Pontificia Universidad Católica del Perú (PUCP), first in my graduating class, with a licentiate thesis approved with distinction. I research electoral behaviour, legislative behaviour and public policy implementation in Peru with mixed methods and R. I am a researcher at the Electoral and Governance School of Peru's National Elections Jury, a consultant to the Corruption Observatory of the Lima Bar Association and a teaching assistant in Political Economy at PUCP.

\seccion{Education}
\begin{entradas}
\entrada[2019 -- 2024]{\cargo{Licentiate in Political Science and Government}{Pontificia Universidad Católica del Perú (PUCP)}
  \detalle{First in the graduating class. Thesis approved with distinction (\textit{mención sobresaliente}, 2025).}}
\entrada[2024]{\cargo{Exchange semester in Political and Social Sciences}{Ghent University, Belgium}
  \detalle{Master's-level courses: Quantitative Methods for Social Research, Public Management, Innovation Management and European Policy.}}
\entrada[2016 -- 2018]{\cargo{International Baccalaureate Diploma}{Colegio de Alto Rendimiento (COAR) Amazonas, Peru}
  \detalle{Graduated with distinction.}}
\end{entradas}

\seccion{Research experience}
\begin{entradas}
\entrada[Jul 2026 --]{\cargo{Researcher}{Electoral and Governance School (ESEG), National Elections Jury of Peru (JNE)}
  \begin{puntos}
  \item I contribute to the design of the institutional research agenda with quantitative and qualitative methods.
  \item I work on the census of electoral studies teaching in Peruvian universities (ongoing).
  \item Technical liaison with academic and legal authorities in electoral studies.
  \end{puntos}}
\entrada[2025 --]{\cargo{Research member}{State and Society Research Group (GIES), PUCP}
  \begin{puntos}
  \item Legislative tracking of four political and institutional reforms ahead of the 2026 general election: I systematised roll-call votes and built the candidates' voting records for the civic portal \enlace{https://hechosynopalabras.pe}{hechosynopalabras.pe}.
  \end{puntos}}
\entrada[Oct 2025 -- \newline Feb 2026]{\cargo{Research assistant}{Centre for Sociological, Economic, Political and Anthropological Research (CISEPA), PUCP}
  \begin{puntos}
  \item I supported the analysis of Peru's political regime (legislative autocratisation and cabinet instability, 2021--2025) for an article submitted to the \textit{Revista de Ciencia Política} (Chile).
  \item I systematised congressional votes from 2021 to 2025 to identify cross-party coalitions on authoritarian issues.
  \item Econometric support in R (interrupted time series, two-way fixed effects and shift-share instruments) for two manuscripts on the effects of the ban on subnational re-election on corruption in public procurement and on municipal fiscal management, submitted to international indexed journals.
  \item I worked on building a corpus of 747,836 headlines from \textit{El Comercio} and \textit{La República} (2016--2024) to measure media coverage of five illicit economies.
  \end{puntos}}
\entrada[Sep 2024 -- \newline Sep 2025]{\cargo{Research assistant}{Q-LAB, Laboratory for AI and Computational Methods in the Social Sciences, PUCP}
  \begin{puntos}
  \item In a project with Swiss Cooperation (GFP Subnacional), I worked on a supervised machine learning model (XGBoost and SHAP values) of delays in 40,243 municipal public works (2018--2024), using records from MEF, RENAMU and INEI; the findings fed into an institutional policy brief.
  \item I built and cleaned district- and region-level datasets for research on corruption, municipal performance and migration.
  \end{puntos}}
\entrada[May -- Aug 2024]{\cargo{Research assistant}{Daniel Encinas (Northwestern University) and Rodrigo Barrenechea}
  \begin{puntos}
  \item I processed province-level results of ten presidential elections (1980--2021) to identify geographic patterns of the populist vote in Peru; descriptive analysis with LAPOP data.
  \end{puntos}}
\entrada[Feb -- Dec 2023]{\cargo{Research intern}{Institute for Democracy and Human Rights (IDEHPUCP)}
  \begin{puntos}
  \item I supported three research projects on governance, migration and human rights, in coordination with United Nations agencies and public-sector counterparts.
  \item I supported the editorial process of the volume marking twenty years of Peru's Truth and Reconciliation Commission Final Report and co-wrote six columns for the IDEHPUCP bulletin.
  \end{puntos}}
\end{entradas}

\seccion{Consulting}
\begin{entradas}
\entrada[Sep 2026 --]{\cargo{Consultant}{Corruption Observatory of the Lima Bar Association}
  \begin{puntos}
  \item I advise on the first instalment of the Observatory's report on public procurement 2016--2025: feasibility and validity of indicators built from open SEACE and OECE records.
  \end{puntos}}
\entrada[Feb -- Aug 2026]{\cargo{Research consultant}{Peruvian School Debate League (LPDE)}
  \begin{puntos}
  \item Applied research on the school debate ecosystem in Peru: study design and quantitative and qualitative analysis.
  \item I prepared the League's first report on that ecosystem (70 institutions) and seven policy briefs for its strategic decisions.
  \end{puntos}}
\entrada[Nov -- Dec 2025]{\cargo{Consultant}{ProDiálogo, commissioned by the National Elections Jury (JNE)}
  \begin{puntos}
  \item I was part of the technical team for three instruments measuring electoral behaviour for the 2026 general election: the Electoral Profile Survey, two survey experiments and a behavioural map.
  \item I took part in segmenting the Peruvian electorate into four profiles, with cluster analysis of surveys from three electoral cycles (2011, 2016 and 2021), as input for the JNE's civic education strategies.
  \end{puntos}}
\entrada[Jul -- Sep 2025]{\cargo{Data analysis consultant}{Sapienza Consultora, for Roche Pharma Peru}
  \begin{puntos}
  \item Analysis of health-sector administrative records (2019--2025) and three reports for the client.
  \end{puntos}}
\end{entradas}

\seccion{Teaching}
\begin{entradas}
\entrada[Aug 2026 --]{\cargo{Teaching assistant}{Political Economy, Faculty of Social Sciences, PUCP}}
\entrada[2024 -- 2025]{\cargo{Teaching assistant}{Q-LAB, PUCP}
  \detalle{Introduction to R for Social Research (two editions) and the Causal Machine Learning module of the Advanced Applied Econometrics course.}}
\end{entradas}

\seccion{Debate leadership and coaching}
\begin{entradas}
\entrada[2026 --]{\cargo{Director General}{Peruvian National Schools Debate Team}
  \detalle{I run the team's call for debaters, coordinate with the heads of the junior and senior squads and represent the team before the organisers of the world schools championship held in Spanish (MED).}}
\entrada[2024 --]{\cargo{Coach}{school and university teams}
  \begin{puntos}
  \item Peruvian National Schools Debate Team (2025--2026): part of the coaching staff for MED Chile 2025 and MED Panama 2026, where the team won its fourth consecutive title; the team I coached reached the quarterfinals and placed among the tournament's top 10 speakers.
  \item Co-head coach of the PUCP Debate Society (since July 2026) and debate coach at Colegio Alpamayo (since June 2026).
  \item University pairs from the Dominican Republic, online (since 2024): in 2025, the first all-Dominican pair to reach the round of 16 at the world universities championship held in Spanish (CMUDE); in 2026, the quarterfinals and the round of 16 (CMUDE Ecuador).
  \end{puntos}}
\entrada[2025 -- 2026]{\cargo{Adjudication}{school and university tournaments in Costa Rica, Ecuador, Peru and Venezuela}
  \detalle{Judge of the CMUDE 2026 grand final (Ecuador) and chief adjudicator of Peru's national schools tournament (TNDE) 2025.}}
\entrada[2022 -- 2025]{\cargo{As a debater}{national inter-university champion (CNDI 2025)}
  \detalle{Semifinalist at CMUDE Madrid 2022, quarterfinalist at CMUDE Panama 2023 and runner-up at the Torneo Rosarista (Bogotá, 2023).}}
\end{entradas}

\seccion{Publications and reports}
\begin{entradas}
\entrada[2026]{Monard, Encinas, Villena-Tagle and \textbf{Valverde} (technical team). \textit{Análisis del elector peruano: aportes desde las ciencias del comportamiento} [The Peruvian voter: contributions from the behavioural sciences]. Lima: National Elections Jury and ProDiálogo, 85~pp.}
\entrada[2025]{\textbf{Valverde, L.} \textit{Desastres y brechas de implementación} [Disasters and implementation gaps: the reconstruction of school infrastructure in Piura, July 2020 -- January 2023, after the coastal El Niño]. Licentiate thesis with distinction, PUCP. \enlace{https://hdl.handle.net/20.500.12404/32370}{PUCP repository}.}
\entrada[2023]{Cabrales, S., Quiroz, S. J., Mamede, T., Barbosa, I. and \textbf{Valverde, L.} (editorial team). \textit{Emergencia climática y derechos humanos} [Climate emergency and human rights: a submission on the advisory opinion before the Inter-American Court of Human Rights]. Alliance of Environmental Legal Clinics of Latin America and the Caribbean.}
\entrada[2023]{Six co-authored columns in the IDEHPUCP bulletin (with I. Jave, M. Cépeda, A. Domínguez, E. Solís and A. Benites), in Spanish, on the search for missing persons, the RENADE registry, teaching about the internal armed conflict, climate change in the courts, disaster prevention and democratic stagnation in Peru.}
\end{entradas}

\seccion{Awards and scholarships}
\begin{entradas}
\entrada[2024]{DARI-PUCP mobility scholarship for the exchange at Ghent University.}
\entrada[2023]{First place, Research Competition of the PUCP Faculty of Social Sciences (undergraduate), with a study of post-disaster reconstruction in Piura.}
\entrada[2019 -- 2024]{Beca 18 (PRONABEC): Peru's national merit scholarship, covering the whole undergraduate degree.}
\end{entradas}

\seccion{Media and panels}
\begin{entradas}
\entrada[2026]{\cargo{Interviews}{\textit{El Salto} (Spain, July), Radio Exitosa (June and July) and Canal 04 Catacaos, on the president-elect, fraud narratives and national politics ahead of the 2026 general election.}}
\entrada[2026]{\cargo{Panels}{\textit{La República}, on the Beca 18 scholarship programme (July), and the Colegio de Politólogos del Perú, on the first presidential debate (March).}}
\end{entradas}

\seccion{Specialised courses}
\begin{entradas}
\entrada[2024 -- 2025]{\cargo{Q-LAB, PUCP}{computational methods}
  \detalle{Recent innovations in regression discontinuity for causal inference (2025); Transformers in Large Language Models for Social Sciences and Public Policy (2025); introduction to machine learning for the social sciences and public management (2024); advanced applied econometrics (2024).}}
\entrada[2023]{\cargo{IDEHPUCP}{human rights}
  \detalle{International humanitarian law, security forces and human rights (advanced training); introduction to environmental management with an intercultural approach.}}
\entrada[2021]{\cargo{Korea Foundation}{The economies of South Korea and Latin America}}
\end{entradas}

\seccion{Methods, tools and languages}
\begin{entradas}
\entrada[Quantitative]{Descriptive and inferential statistics; econometrics and causal inference; applied machine learning (XGBoost, SHAP); cluster analysis; spatial analysis.}
\entrada[Qualitative]{Case studies; process tracing; semi-structured interviews; document, content and discourse analysis.}
\entrada[Political analysis]{Legislative and regulatory tracking; institutional and current-affairs analysis; stakeholder mapping.}
\entrada[Software]{R, my main tool (tidyverse, fixest, tidymodels); Python (scikit-learn, spaCy) and Stata at an intermediate level; Git and GitHub; LaTeX; KoboToolbox.}
\entrada[Languages]{Spanish, native. English: Cambridge First Certificate in English (B2); master's-level courses taught in English at Ghent University.}
\end{entradas}

\seccion{References}
\ifdefined\conreferencias
  \input{privado/referencias-en}
\else
  Available on request.
\fi

\end{document}
